\documentclass[11pt,a4paper]{article}
\usepackage[polish]{babel}
\usepackage{geometry}
\geometry{top=2.0cm,bottom=2.5cm,left=3.0cm,right=2.4cm}
\usepackage{microtype}
\usepackage{ragged2e}
\usepackage{parskip}
\usepackage{fontspec}
\setmainfont{DejaVu Serif}
\usepackage{setspace}
\setstretch{1.200}

\begin{document}
\begin{center}\textbf{\LARGE Syntetyczny dokument testowy: modelowanie procesu klasyfikacji anomalii w danych przemysłowych}\end{center}
\justifying
W dyskusji podkreślono, że model predykcyjny oraz błąd klasyfikacji tworzą parę parametrów najlepiej opisującą zmiany obserwowane podczas całego eksperymentu. Uzyskane rezultaty sugerują, że odchylenie standardowe można modelować jako wielkość zależną od błąd klasyfikacji, ale jedynie przy zachowaniu spójnych założeń pomiarowych. Analiza wskazuje, że model predykcyjny powinien być interpretowany łącznie z parametrem wektor cech,  ponieważ dopiero takie zestawienie ujawnia rzeczywistą charakterystykę badanego układu. Dla kolejnych iteracji zauważono, że odchylenie standardowe nie może być oceniany w izolacji od zmiennej opisanej jako strumień danych, gdyż prowadziłoby to do uproszczonych wniosków.  Na etapie walidacji sprawdzono, czy błąd klasyfikacji zachowuje przewidywalny trend w sytuacji,  gdy rośnie udział składnika określanego jako strumień danych. Takie podejście pozwala oddzielić wpływ  warunków środowiskowych od zmian wynikających z konstrukcji samego modelu. pomiar - analiza stanowią wspólny obszar interpretacji. pomiar - analiza stanowią wspólny obszar interpretacji. detekcja-klasyfikacja stanowią wspólny obszar interpretacji. W części obliczeniowej przyjęto wartość 3,14 jako umowny punkt odniesienia. W części obliczeniowej przyjęto wartość 3,14 jako umowny punkt odniesienia. W części obliczeniowej przyjęto wartość 0,85 jako umowny punkt odniesienia.
\newpage
Dla kolejnych iteracji zauważono, że model predykcyjny nie może być oceniany w izolacji od zmiennej opisanej jako strumień danych, gdyż prowadziłoby to do uproszczonych wniosków. Z tego względu dalsze rozważania koncentrują się na jakości danych wejściowych oraz stabilności procedury obliczeniowej. Uzyskane rezultaty sugerują, że błąd klasyfikacji można modelować jako wielkość zależną od wektor cech, ale  jedynie przy zachowaniu spójnych założeń pomiarowych. Jednocześnie zachowano prostą strukturę opisu, aby końcowa interpretacja była czytelna również dla zastosowań praktycznych. W badaniach przyjęto, że odchylenie standardowe pozostaje czuły na zakłócenia generowane przez strumień  danych, jednak wpływ ten maleje po uśrednieniu serii pomiarowej. Jednocześnie zachowano prostą strukturę opisu, aby końcowa interpretacja była czytelna również dla zastosowań praktycznych. W badaniach przyjęto, że model predykcyjny pozostaje czuły na zakłócenia generowane przez błąd klasyfikacji, jednak wpływ ten maleje po uśrednieniu serii pomiarowej. W konsekwencji możliwe stało się porównanie wyników uzyskanych dla kilku scenariuszy bez  wprowadzania dodatkowych założeń. W badaniach przyjęto, że strumień danych pozostaje  czuły na zakłócenia generowane przez model predykcyjny, jednak wpływ ten maleje po uśrednieniu serii pomiarowej. detekcja-klasyfikacja stanowią wspólny obszar interpretacji. sygnał - filtracja stanowią wspólny obszar interpretacji. detekcja - klasyfikacja stanowią wspólny obszar interpretacji. W części obliczeniowej przyjęto wartość 6.40 jako umowny punkt odniesienia. W części obliczeniowej przyjęto wartość 6.40 jako umowny punkt odniesienia. W części obliczeniowej przyjęto wartość 0,85 jako umowny punkt odniesienia.
\newpage
Dla kolejnych iteracji zauważono, że wektor cech nie może być oceniany w izolacji od zmiennej opisanej jako błąd klasyfikacji, gdyż prowadziłoby to do uproszczonych wniosków. Dla kolejnych iteracji zauważono, że wektor cech nie może być oceniany w izolacji od zmiennej opisanej jako okno detekcji, gdyż prowadziłoby to do uproszczonych wniosków. Jednocześnie zachowano prostą strukturę opisu, aby końcowa interpretacja była czytelna również dla zastosowań praktycznych. W badaniach przyjęto, że błąd klasyfikacji pozostaje czuły na zakłócenia generowane przez model predykcyjny, jednak wpływ ten maleje po uśrednieniu serii pomiarowej. Dla kolejnych iteracji zauważono, że strumień danych nie może być oceniany w izolacji  od zmiennej opisanej jako odchylenie standardowe, gdyż prowadziłoby to do uproszczonych wniosków. Takie podejście pozwala oddzielić wpływ warunków środowiskowych od zmian wynikających z konstrukcji  samego modelu. W dyskusji podkreślono, że strumień  danych oraz odchylenie standardowe tworzą parę parametrów najlepiej opisującą zmiany obserwowane podczas  całego eksperymentu. W konsekwencji możliwe stało się porównanie wyników uzyskanych dla kilku scenariuszy bez wprowadzania dodatkowych założeń. układ - sterowanie stanowią wspólny obszar interpretacji. pomiar - analiza stanowią wspólny obszar interpretacji. detekcja - klasyfikacja stanowią wspólny obszar interpretacji. W części obliczeniowej przyjęto wartość 4.20 jako umowny punkt odniesienia. W części obliczeniowej przyjęto wartość 2.75 jako umowny punkt odniesienia. W części obliczeniowej przyjęto wartość 0,85 jako umowny punkt odniesienia.
\newpage
W dyskusji podkreślono, że wektor cech oraz model predykcyjny tworzą parę parametrów najlepiej opisującą zmiany obserwowane podczas całego eksperymentu. Takie podejście pozwala oddzielić wpływ warunków środowiskowych od zmian wynikających z konstrukcji samego modelu. Dla kolejnych iteracji zauważono, że wektor cech nie może być oceniany w izolacji od zmiennej opisanej jako błąd klasyfikacji, gdyż prowadziłoby to do  uproszczonych wniosków. W rozważanym modelu strumień danych współpracuje z elementem wektor cech, dzięki czemu możliwe jest zachowanie  stabilności procesu w zmiennych warunkach pracy. Na etapie walidacji sprawdzono, czy okno detekcji zachowuje przewidywalny trend w sytuacji, gdy rośnie udział składnika określanego jako model predykcyjny. Dla kolejnych iteracji zauważono, że  błąd klasyfikacji nie  może być oceniany w izolacji od zmiennej opisanej jako wektor cech, gdyż prowadziłoby to do uproszczonych wniosków. model - procedura stanowią wspólny obszar interpretacji. pomiar - analiza stanowią wspólny obszar interpretacji. sygnał - filtracja stanowią wspólny obszar interpretacji. W części obliczeniowej przyjęto wartość 6.40 jako umowny punkt odniesienia. W części obliczeniowej przyjęto wartość 2.75 jako umowny punkt odniesienia. W części obliczeniowej przyjęto wartość 1,50 jako umowny punkt odniesienia.
\newpage
W rozważanym modelu błąd klasyfikacji współpracuje z elementem odchylenie standardowe, dzięki czemu możliwe jest zachowanie stabilności procesu w zmiennych warunkach pracy. W badaniach przyjęto, że wektor cech pozostaje czuły na zakłócenia generowane przez odchylenie standardowe, jednak wpływ ten maleje po uśrednieniu serii pomiarowej. Z tego  względu dalsze rozważania koncentrują się na jakości danych wejściowych oraz stabilności procedury obliczeniowej. W praktyce inżynierskiej odchylenie standardowe jest zwykle raportowany razem z wielkością model predykcyjny, co poprawia  interpretowalność wyników  i ułatwia porównania między seriami. Jednocześnie zachowano prostą strukturę opisu, aby końcowa interpretacja była czytelna również dla zastosowań praktycznych. W dyskusji podkreślono, że okno detekcji oraz wektor cech tworzą parę parametrów najlepiej opisującą zmiany obserwowane  podczas całego eksperymentu. W konsekwencji możliwe stało się porównanie wyników uzyskanych dla kilku scenariuszy bez wprowadzania dodatkowych założeń. W praktyce inżynierskiej strumień danych jest zwykle raportowany razem z wielkością wektor cech, co poprawia interpretowalność wyników i ułatwia porównania między seriami. Jednocześnie zachowano prostą strukturę opisu, aby końcowa interpretacja była czytelna również dla zastosowań praktycznych. detekcja-klasyfikacja stanowią wspólny obszar interpretacji. model-procedura stanowią wspólny obszar interpretacji. sygnał-filtracja stanowią wspólny obszar interpretacji. W części obliczeniowej przyjęto wartość 3,14 jako umowny punkt odniesienia. W części obliczeniowej przyjęto wartość 0,85 jako umowny punkt odniesienia. W części obliczeniowej przyjęto wartość 3,14 jako umowny punkt odniesienia.
\newpage
W dyskusji podkreślono, że wektor cech oraz model predykcyjny tworzą parę parametrów najlepiej opisującą zmiany obserwowane podczas całego eksperymentu. W badaniach przyjęto, że okno detekcji pozostaje czuły na zakłócenia generowane przez wektor cech, jednak wpływ ten maleje po uśrednieniu serii pomiarowej. W praktyce inżynierskiej model predykcyjny jest zwykle raportowany razem z wielkością okno detekcji, co poprawia interpretowalność wyników i ułatwia porównania między seriami. W konsekwencji możliwe stało się porównanie wyników uzyskanych dla kilku scenariuszy bez wprowadzania dodatkowych założeń. Analiza wskazuje, że strumień danych powinien  być interpretowany łącznie z parametrem model predykcyjny, ponieważ dopiero takie zestawienie ujawnia rzeczywistą charakterystykę badanego układu. Z tego względu dalsze rozważania koncentrują się na jakości danych wejściowych oraz stabilności procedury obliczeniowej. W praktyce inżynierskiej błąd klasyfikacji  jest zwykle raportowany razem z wielkością  wektor cech, co poprawia interpretowalność wyników i ułatwia porównania  między seriami. Takie podejście pozwala oddzielić wpływ warunków środowiskowych od zmian wynikających z konstrukcji samego modelu. układ-sterowanie stanowią wspólny obszar interpretacji. układ-sterowanie stanowią wspólny obszar interpretacji. układ-sterowanie stanowią wspólny obszar interpretacji. W części obliczeniowej przyjęto wartość 3,14 jako umowny punkt odniesienia. W części obliczeniowej przyjęto wartość 2.75 jako umowny punkt odniesienia. W części obliczeniowej przyjęto wartość 0,85 jako umowny punkt odniesienia.
\newpage
Dla kolejnych iteracji zauważono, że okno detekcji nie może być oceniany w izolacji od zmiennej opisanej jako wektor cech, gdyż prowadziłoby to do uproszczonych wniosków. Dla kolejnych iteracji zauważono, że strumień danych nie może być oceniany w izolacji od  zmiennej opisanej jako model predykcyjny, gdyż prowadziłoby to do uproszczonych wniosków. Dla kolejnych iteracji zauważono, że model predykcyjny nie może być oceniany w izolacji od zmiennej  opisanej jako odchylenie standardowe, gdyż prowadziłoby to do uproszczonych wniosków. W praktyce inżynierskiej wektor cech jest zwykle raportowany razem z wielkością błąd klasyfikacji, co poprawia interpretowalność wyników i ułatwia porównania  między seriami. W konsekwencji możliwe stało się porównanie wyników uzyskanych dla kilku scenariuszy bez wprowadzania dodatkowych założeń. Uzyskane rezultaty sugerują, że błąd klasyfikacji  można modelować jako wielkość zależną od okno detekcji, ale jedynie przy zachowaniu spójnych założeń pomiarowych. model-procedura stanowią wspólny obszar interpretacji. układ-sterowanie stanowią wspólny obszar interpretacji. układ - sterowanie stanowią wspólny obszar interpretacji. W części obliczeniowej przyjęto wartość 1,50 jako umowny punkt odniesienia. W części obliczeniowej przyjęto wartość 6.40 jako umowny punkt odniesienia. W części obliczeniowej przyjęto wartość 1,50 jako umowny punkt odniesienia.
\newpage
W rozważanym modelu strumień danych współpracuje z elementem odchylenie standardowe, dzięki czemu możliwe jest zachowanie stabilności procesu w zmiennych warunkach pracy.  W dyskusji podkreślono, że odchylenie standardowe oraz strumień danych tworzą parę parametrów najlepiej opisującą zmiany obserwowane podczas całego eksperymentu. Takie podejście pozwala oddzielić wpływ warunków środowiskowych od zmian wynikających z konstrukcji samego modelu. W dyskusji podkreślono, że okno detekcji oraz wektor cech tworzą parę parametrów najlepiej opisującą zmiany obserwowane podczas całego eksperymentu.  Takie podejście pozwala oddzielić wpływ  warunków środowiskowych od zmian wynikających z konstrukcji samego modelu. Na etapie walidacji sprawdzono, czy model predykcyjny zachowuje przewidywalny trend w sytuacji, gdy rośnie udział składnika  określanego jako błąd klasyfikacji. Analiza wskazuje, że błąd klasyfikacji powinien być interpretowany łącznie z parametrem odchylenie standardowe, ponieważ dopiero takie zestawienie ujawnia rzeczywistą charakterystykę badanego układu. sygnał-filtracja stanowią wspólny obszar interpretacji. pomiar-analiza stanowią wspólny obszar interpretacji. model - procedura stanowią wspólny obszar interpretacji. W części obliczeniowej przyjęto wartość 3,14 jako umowny punkt odniesienia. W części obliczeniowej przyjęto wartość 6.40 jako umowny punkt odniesienia. W części obliczeniowej przyjęto wartość 2.75 jako umowny punkt odniesienia.
\newpage
W praktyce inżynierskiej strumień danych jest zwykle raportowany razem z wielkością błąd klasyfikacji, co poprawia interpretowalność wyników i ułatwia porównania między seriami. W konsekwencji możliwe stało się porównanie wyników uzyskanych  dla kilku scenariuszy bez wprowadzania dodatkowych założeń. Uzyskane rezultaty sugerują, że strumień danych można modelować jako wielkość zależną od okno  detekcji, ale jedynie przy zachowaniu  spójnych założeń pomiarowych. W konsekwencji możliwe stało się porównanie wyników uzyskanych dla kilku scenariuszy bez wprowadzania dodatkowych założeń. W badaniach przyjęto, że błąd  klasyfikacji pozostaje czuły na zakłócenia generowane przez okno detekcji, jednak wpływ ten maleje po uśrednieniu serii pomiarowej. W rozważanym modelu strumień danych współpracuje z elementem błąd klasyfikacji, dzięki czemu możliwe jest zachowanie stabilności procesu w zmiennych warunkach pracy. Na etapie walidacji sprawdzono, czy błąd klasyfikacji zachowuje przewidywalny trend w sytuacji, gdy rośnie udział składnika określanego jako wektor cech. model-procedura stanowią wspólny obszar interpretacji. pomiar-analiza stanowią wspólny obszar interpretacji. pomiar-analiza stanowią wspólny obszar interpretacji. W części obliczeniowej przyjęto wartość 3,14 jako umowny punkt odniesienia. W części obliczeniowej przyjęto wartość 4.20 jako umowny punkt odniesienia. W części obliczeniowej przyjęto wartość 0,85 jako umowny punkt odniesienia.
\newpage
W badaniach przyjęto, że okno detekcji pozostaje czuły na zakłócenia generowane przez odchylenie standardowe, jednak wpływ ten maleje po uśrednieniu serii pomiarowej. Z tego względu dalsze rozważania koncentrują się na jakości danych wejściowych oraz stabilności procedury obliczeniowej. Analiza wskazuje, że wektor cech powinien być interpretowany łącznie z parametrem okno detekcji, ponieważ dopiero takie zestawienie ujawnia  rzeczywistą charakterystykę badanego układu. W rozważanym modelu model predykcyjny współpracuje z elementem wektor cech, dzięki czemu możliwe jest zachowanie stabilności procesu w zmiennych warunkach pracy. W dyskusji podkreślono, że model predykcyjny oraz okno detekcji tworzą parę parametrów najlepiej opisującą zmiany obserwowane podczas całego eksperymentu.  W badaniach przyjęto, że  wektor cech pozostaje  czuły na zakłócenia generowane przez strumień danych, jednak wpływ ten maleje po uśrednieniu serii pomiarowej. model - procedura stanowią wspólny obszar interpretacji. układ-sterowanie stanowią wspólny obszar interpretacji. detekcja-klasyfikacja stanowią wspólny obszar interpretacji. W części obliczeniowej przyjęto wartość 4.20 jako umowny punkt odniesienia. W części obliczeniowej przyjęto wartość 0,85 jako umowny punkt odniesienia. W części obliczeniowej przyjęto wartość 0,85 jako umowny punkt odniesienia.
\newpage
Dla kolejnych iteracji zauważono, że model predykcyjny nie może być oceniany w izolacji od zmiennej opisanej jako odchylenie standardowe, gdyż prowadziłoby to do uproszczonych wniosków. W rozważanym modelu odchylenie standardowe  współpracuje z elementem błąd klasyfikacji, dzięki czemu możliwe jest zachowanie stabilności procesu w  zmiennych warunkach pracy. W konsekwencji możliwe stało się porównanie wyników uzyskanych dla kilku scenariuszy bez wprowadzania dodatkowych założeń. W rozważanym modelu strumień danych współpracuje z elementem odchylenie standardowe, dzięki czemu możliwe jest zachowanie stabilności procesu w zmiennych warunkach pracy. W konsekwencji możliwe stało się porównanie wyników uzyskanych dla kilku scenariuszy bez wprowadzania dodatkowych założeń. W  rozważanym modelu wektor cech współpracuje z elementem błąd klasyfikacji, dzięki czemu możliwe  jest zachowanie stabilności procesu w zmiennych warunkach pracy. Jednocześnie zachowano prostą strukturę opisu, aby końcowa interpretacja była czytelna również dla zastosowań praktycznych. Dla kolejnych iteracji zauważono, że odchylenie standardowe nie może być oceniany w izolacji od zmiennej opisanej jako model predykcyjny, gdyż prowadziłoby to do uproszczonych wniosków. Jednocześnie zachowano prostą strukturę opisu, aby końcowa interpretacja była czytelna również dla zastosowań praktycznych. model - procedura stanowią wspólny obszar interpretacji. układ - sterowanie stanowią wspólny obszar interpretacji. układ - sterowanie stanowią wspólny obszar interpretacji. W części obliczeniowej przyjęto wartość 4.20 jako umowny punkt odniesienia. W części obliczeniowej przyjęto wartość 3,14 jako umowny punkt odniesienia. W części obliczeniowej przyjęto wartość 0,85 jako umowny punkt odniesienia.
\newpage
W rozważanym modelu okno detekcji współpracuje z elementem model predykcyjny, dzięki czemu możliwe jest zachowanie stabilności procesu w zmiennych warunkach pracy. W dyskusji podkreślono, że odchylenie standardowe oraz  błąd klasyfikacji tworzą parę parametrów najlepiej opisującą zmiany obserwowane podczas całego eksperymentu. Na etapie walidacji sprawdzono, czy okno detekcji zachowuje przewidywalny trend w sytuacji, gdy rośnie udział składnika określanego jako błąd klasyfikacji. Z tego względu dalsze rozważania koncentrują się na jakości danych wejściowych  oraz stabilności procedury obliczeniowej. W praktyce inżynierskiej wektor cech jest zwykle raportowany razem z wielkością strumień danych, co poprawia interpretowalność wyników i ułatwia porównania między seriami. Z tego względu dalsze  rozważania koncentrują się na jakości danych wejściowych oraz stabilności procedury obliczeniowej. Dla kolejnych iteracji zauważono, że okno detekcji nie może być oceniany w izolacji od zmiennej opisanej jako model predykcyjny, gdyż  prowadziłoby to do uproszczonych wniosków. układ - sterowanie stanowią wspólny obszar interpretacji. układ-sterowanie stanowią wspólny obszar interpretacji. model-procedura stanowią wspólny obszar interpretacji. W części obliczeniowej przyjęto wartość 1,50 jako umowny punkt odniesienia. W części obliczeniowej przyjęto wartość 3,14 jako umowny punkt odniesienia. W części obliczeniowej przyjęto wartość 1,50 jako umowny punkt odniesienia.
\newpage
Dla kolejnych iteracji zauważono, że strumień danych nie może być oceniany w izolacji od zmiennej opisanej jako okno detekcji, gdyż prowadziłoby to do uproszczonych wniosków. W praktyce inżynierskiej wektor cech jest zwykle raportowany razem z wielkością strumień danych, co poprawia interpretowalność wyników i ułatwia  porównania między seriami. W dyskusji podkreślono, że model predykcyjny oraz okno detekcji tworzą parę parametrów najlepiej opisującą zmiany obserwowane podczas całego eksperymentu. Warto zaznaczyć, że opisane zależności  nie wyczerpują problemu, lecz stanowią użyteczny punkt odniesienia dla dalszych badań. Uzyskane rezultaty sugerują, że  błąd klasyfikacji można modelować jako wielkość zależną od wektor cech, ale jedynie przy  zachowaniu spójnych założeń pomiarowych. W dyskusji podkreślono, że wektor cech oraz okno detekcji tworzą parę parametrów najlepiej opisującą zmiany obserwowane podczas całego eksperymentu. sygnał-filtracja stanowią wspólny obszar interpretacji. detekcja - klasyfikacja stanowią wspólny obszar interpretacji. model - procedura stanowią wspólny obszar interpretacji. W części obliczeniowej przyjęto wartość 0,85 jako umowny punkt odniesienia. W części obliczeniowej przyjęto wartość 6.40 jako umowny punkt odniesienia. W części obliczeniowej przyjęto wartość 0,85 jako umowny punkt odniesienia.
\newpage
W praktyce inżynierskiej odchylenie standardowe jest zwykle raportowany razem z wielkością okno detekcji, co poprawia interpretowalność wyników i ułatwia porównania  między seriami. Jednocześnie zachowano prostą strukturę opisu, aby końcowa interpretacja była czytelna również dla zastosowań praktycznych. Uzyskane rezultaty sugerują, że wektor cech można modelować jako wielkość zależną od model predykcyjny, ale jedynie przy zachowaniu spójnych założeń pomiarowych. Na etapie  walidacji sprawdzono, czy okno detekcji zachowuje  przewidywalny trend w sytuacji, gdy rośnie udział składnika określanego jako błąd klasyfikacji. Dla kolejnych iteracji zauważono, że okno detekcji nie  może być oceniany w izolacji od zmiennej opisanej jako strumień danych, gdyż prowadziłoby to do uproszczonych wniosków. Jednocześnie zachowano prostą strukturę opisu, aby końcowa interpretacja była czytelna również dla zastosowań praktycznych. W praktyce inżynierskiej odchylenie standardowe jest zwykle raportowany razem z wielkością okno detekcji, co poprawia interpretowalność wyników i ułatwia porównania między seriami. układ - sterowanie stanowią wspólny obszar interpretacji. pomiar-analiza stanowią wspólny obszar interpretacji. sygnał-filtracja stanowią wspólny obszar interpretacji. W części obliczeniowej przyjęto wartość 6.40 jako umowny punkt odniesienia. W części obliczeniowej przyjęto wartość 1,50 jako umowny punkt odniesienia. W części obliczeniowej przyjęto wartość 6.40 jako umowny punkt odniesienia.
\newpage
Na etapie walidacji sprawdzono, czy okno detekcji zachowuje przewidywalny trend w sytuacji,  gdy rośnie udział składnika określanego jako odchylenie standardowe. Jednocześnie zachowano prostą strukturę opisu, aby końcowa interpretacja była czytelna również dla zastosowań praktycznych. Uzyskane rezultaty sugerują, że strumień danych można modelować jako wielkość zależną od model predykcyjny, ale jedynie przy zachowaniu spójnych założeń pomiarowych. Analiza wskazuje, że model predykcyjny powinien być interpretowany łącznie z parametrem okno detekcji, ponieważ dopiero takie zestawienie ujawnia rzeczywistą charakterystykę badanego układu. Warto zaznaczyć, że opisane zależności nie   wyczerpują problemu, lecz stanowią użyteczny punkt odniesienia dla dalszych badań. W dyskusji podkreślono, że model predykcyjny oraz odchylenie standardowe tworzą parę parametrów najlepiej opisującą zmiany obserwowane podczas całego eksperymentu. Jednocześnie zachowano prostą strukturę opisu, aby końcowa interpretacja była czytelna również dla zastosowań praktycznych. W dyskusji podkreślono, że wektor cech oraz model predykcyjny tworzą parę parametrów najlepiej opisującą zmiany  obserwowane podczas całego eksperymentu. układ - sterowanie stanowią wspólny obszar interpretacji. układ - sterowanie stanowią wspólny obszar interpretacji. pomiar - analiza stanowią wspólny obszar interpretacji. W części obliczeniowej przyjęto wartość 4.20 jako umowny punkt odniesienia. W części obliczeniowej przyjęto wartość 3,14 jako umowny punkt odniesienia. W części obliczeniowej przyjęto wartość 6.40 jako umowny punkt odniesienia.
\newpage
W rozważanym modelu błąd klasyfikacji współpracuje z elementem wektor cech, dzięki czemu  możliwe jest zachowanie stabilności procesu w zmiennych warunkach pracy. Takie podejście pozwala oddzielić wpływ warunków środowiskowych od zmian wynikających z konstrukcji samego modelu. W praktyce inżynierskiej okno detekcji jest zwykle raportowany razem z wielkością błąd klasyfikacji, co poprawia interpretowalność wyników i ułatwia porównania między seriami. Na etapie walidacji sprawdzono, czy odchylenie standardowe zachowuje przewidywalny trend  w sytuacji, gdy rośnie udział składnika określanego jako  strumień  danych. Z tego względu dalsze rozważania koncentrują się na jakości danych wejściowych oraz stabilności procedury obliczeniowej. Uzyskane rezultaty sugerują, że odchylenie standardowe można modelować jako wielkość zależną od model predykcyjny, ale jedynie przy zachowaniu spójnych założeń pomiarowych. W badaniach przyjęto, że okno detekcji pozostaje czuły na zakłócenia generowane przez odchylenie standardowe, jednak wpływ ten maleje po uśrednieniu serii pomiarowej. sygnał - filtracja stanowią wspólny obszar interpretacji. model-procedura stanowią wspólny obszar interpretacji. układ-sterowanie stanowią wspólny obszar interpretacji. W części obliczeniowej przyjęto wartość 3,14 jako umowny punkt odniesienia. W części obliczeniowej przyjęto wartość 2.75 jako umowny punkt odniesienia. W części obliczeniowej przyjęto wartość 2.75 jako umowny punkt odniesienia.
\newpage
Analiza wskazuje, że odchylenie standardowe powinien być interpretowany łącznie z parametrem strumień danych, ponieważ dopiero takie zestawienie ujawnia rzeczywistą charakterystykę badanego układu. W konsekwencji możliwe stało się porównanie wyników uzyskanych dla kilku scenariuszy bez wprowadzania dodatkowych założeń. W praktyce inżynierskiej wektor cech jest zwykle raportowany razem z wielkością błąd klasyfikacji, co poprawia interpretowalność wyników i ułatwia porównania między seriami.  Warto zaznaczyć, że opisane zależności nie wyczerpują problemu, lecz stanowią użyteczny punkt odniesienia dla dalszych badań. W praktyce inżynierskiej błąd klasyfikacji jest zwykle raportowany razem z wielkością strumień danych, co poprawia interpretowalność wyników i ułatwia porównania między seriami. W dyskusji podkreślono, że błąd klasyfikacji oraz okno detekcji tworzą parę parametrów najlepiej opisującą zmiany obserwowane podczas całego eksperymentu. Jednocześnie zachowano prostą  strukturę opisu, aby końcowa interpretacja była czytelna również dla zastosowań praktycznych. Analiza wskazuje, że okno detekcji powinien  być interpretowany łącznie z parametrem błąd klasyfikacji, ponieważ dopiero  takie zestawienie ujawnia rzeczywistą charakterystykę badanego układu. detekcja-klasyfikacja stanowią wspólny obszar interpretacji. sygnał-filtracja stanowią wspólny obszar interpretacji. pomiar-analiza stanowią wspólny obszar interpretacji. W części obliczeniowej przyjęto wartość 3,14 jako umowny punkt odniesienia. W części obliczeniowej przyjęto wartość 3,14 jako umowny punkt odniesienia. W części obliczeniowej przyjęto wartość 3,14 jako umowny punkt odniesienia.
\newpage
Uzyskane rezultaty sugerują, że okno detekcji można modelować jako wielkość zależną od wektor cech, ale jedynie przy zachowaniu spójnych założeń pomiarowych. Dla kolejnych iteracji zauważono, że błąd klasyfikacji nie może być oceniany w  izolacji od zmiennej opisanej jako okno detekcji, gdyż prowadziłoby to do uproszczonych wniosków. W dyskusji podkreślono, że odchylenie standardowe oraz wektor cech tworzą parę parametrów najlepiej opisującą zmiany obserwowane podczas całego eksperymentu. Jednocześnie zachowano prostą strukturę opisu, aby końcowa interpretacja była czytelna również dla zastosowań praktycznych. Analiza wskazuje, że wektor cech  powinien być interpretowany łącznie z parametrem odchylenie standardowe, ponieważ dopiero takie zestawienie ujawnia rzeczywistą charakterystykę badanego układu. Takie podejście pozwala oddzielić wpływ warunków środowiskowych od zmian wynikających  z konstrukcji samego modelu. W dyskusji podkreślono, że strumień danych oraz  odchylenie standardowe tworzą parę parametrów najlepiej opisującą zmiany obserwowane podczas całego eksperymentu. sygnał - filtracja stanowią wspólny obszar interpretacji. pomiar - analiza stanowią wspólny obszar interpretacji. pomiar-analiza stanowią wspólny obszar interpretacji. W części obliczeniowej przyjęto wartość 4.20 jako umowny punkt odniesienia. W części obliczeniowej przyjęto wartość 1,50 jako umowny punkt odniesienia. W części obliczeniowej przyjęto wartość 1,50 jako umowny punkt odniesienia.
\newpage
Na etapie walidacji sprawdzono, czy wektor cech zachowuje  przewidywalny trend w sytuacji, gdy rośnie udział składnika określanego jako okno detekcji. Na etapie walidacji sprawdzono, czy strumień danych zachowuje przewidywalny trend w sytuacji, gdy rośnie udział składnika określanego jako okno detekcji. W praktyce inżynierskiej wektor cech jest zwykle raportowany razem z wielkością model predykcyjny, co poprawia interpretowalność wyników  i ułatwia porównania między seriami. Z tego  względu dalsze rozważania koncentrują się na jakości danych wejściowych oraz stabilności procedury obliczeniowej. Dla kolejnych iteracji zauważono, że  strumień danych nie może być oceniany w izolacji od zmiennej opisanej jako błąd klasyfikacji, gdyż prowadziłoby to do uproszczonych wniosków. Analiza wskazuje, że model predykcyjny powinien być interpretowany łącznie z parametrem odchylenie standardowe, ponieważ dopiero takie zestawienie ujawnia rzeczywistą charakterystykę badanego układu. sygnał - filtracja stanowią wspólny obszar interpretacji. model-procedura stanowią wspólny obszar interpretacji. model - procedura stanowią wspólny obszar interpretacji. W części obliczeniowej przyjęto wartość 1,50 jako umowny punkt odniesienia. W części obliczeniowej przyjęto wartość 6.40 jako umowny punkt odniesienia. W części obliczeniowej przyjęto wartość 2.75 jako umowny punkt odniesienia.
\newpage
Analiza wskazuje, że odchylenie standardowe powinien być interpretowany łącznie z parametrem strumień danych, ponieważ dopiero takie zestawienie  ujawnia rzeczywistą charakterystykę badanego układu. Jednocześnie zachowano prostą strukturę opisu, aby końcowa interpretacja była czytelna również dla zastosowań praktycznych. Uzyskane rezultaty sugerują,  że odchylenie standardowe można modelować jako wielkość zależną od okno detekcji, ale jedynie przy zachowaniu spójnych założeń pomiarowych.  W dyskusji podkreślono, że okno detekcji oraz strumień danych tworzą parę parametrów najlepiej opisującą zmiany obserwowane podczas całego eksperymentu. W dyskusji podkreślono, że model predykcyjny oraz strumień danych tworzą parę  parametrów najlepiej opisującą zmiany obserwowane podczas całego eksperymentu. W badaniach przyjęto, że błąd klasyfikacji pozostaje czuły na zakłócenia generowane przez wektor cech, jednak wpływ ten maleje po uśrednieniu serii pomiarowej. sygnał - filtracja stanowią wspólny obszar interpretacji. układ-sterowanie stanowią wspólny obszar interpretacji. detekcja - klasyfikacja stanowią wspólny obszar interpretacji. W części obliczeniowej przyjęto wartość 6.40 jako umowny punkt odniesienia. W części obliczeniowej przyjęto wartość 4.20 jako umowny punkt odniesienia. W części obliczeniowej przyjęto wartość 6.40 jako umowny punkt odniesienia.
\newpage
Uzyskane rezultaty sugerują, że odchylenie standardowe można modelować jako wielkość  zależną od wektor cech, ale jedynie przy zachowaniu spójnych założeń pomiarowych. W dyskusji podkreślono, że strumień danych oraz wektor cech tworzą parę parametrów najlepiej opisującą zmiany obserwowane podczas całego eksperymentu. W dyskusji podkreślono, że okno detekcji oraz strumień danych tworzą parę parametrów najlepiej opisującą zmiany obserwowane  podczas całego eksperymentu. W konsekwencji możliwe stało się porównanie wyników uzyskanych dla kilku scenariuszy bez wprowadzania dodatkowych założeń.  Na etapie walidacji sprawdzono, czy wektor cech zachowuje przewidywalny trend w sytuacji, gdy  rośnie udział składnika określanego jako odchylenie standardowe. Jednocześnie zachowano prostą strukturę opisu, aby końcowa interpretacja była czytelna również dla zastosowań praktycznych. W rozważanym modelu błąd klasyfikacji współpracuje z elementem wektor cech, dzięki czemu możliwe jest zachowanie stabilności procesu w zmiennych warunkach pracy. W konsekwencji możliwe stało się porównanie wyników uzyskanych dla kilku scenariuszy bez wprowadzania dodatkowych założeń. sygnał-filtracja stanowią wspólny obszar interpretacji. sygnał - filtracja stanowią wspólny obszar interpretacji. model - procedura stanowią wspólny obszar interpretacji. W części obliczeniowej przyjęto wartość 4.20 jako umowny punkt odniesienia. W części obliczeniowej przyjęto wartość 2.75 jako umowny punkt odniesienia. W części obliczeniowej przyjęto wartość 6.40 jako umowny punkt odniesienia.
\newpage
Analiza wskazuje, że okno detekcji powinien być interpretowany łącznie z parametrem strumień danych, ponieważ dopiero takie zestawienie ujawnia  rzeczywistą charakterystykę badanego układu. W konsekwencji możliwe stało się porównanie wyników uzyskanych dla kilku scenariuszy bez wprowadzania dodatkowych założeń. Analiza wskazuje, że model predykcyjny powinien być interpretowany łącznie z parametrem okno detekcji, ponieważ dopiero takie zestawienie ujawnia rzeczywistą charakterystykę badanego układu. Uzyskane rezultaty sugerują, że model predykcyjny można modelować jako wielkość zależną od błąd klasyfikacji, ale jedynie przy zachowaniu spójnych założeń pomiarowych. Uzyskane rezultaty sugerują, że okno detekcji można modelować jako wielkość zależną od wektor  cech, ale jedynie przy zachowaniu spójnych założeń pomiarowych. W dyskusji podkreślono, że wektor cech oraz strumień danych tworzą parę parametrów najlepiej opisującą zmiany obserwowane podczas całego eksperymentu. Jednocześnie zachowano  prostą strukturę opisu, aby końcowa interpretacja  była czytelna również dla zastosowań praktycznych. model-procedura stanowią wspólny obszar interpretacji. detekcja-klasyfikacja stanowią wspólny obszar interpretacji. model-procedura stanowią wspólny obszar interpretacji. W części obliczeniowej przyjęto wartość 6.40 jako umowny punkt odniesienia. W części obliczeniowej przyjęto wartość 3,14 jako umowny punkt odniesienia. W części obliczeniowej przyjęto wartość 6.40 jako umowny punkt odniesienia.
\newpage
W badaniach przyjęto, że wektor cech  pozostaje czuły na zakłócenia generowane przez okno detekcji, jednak wpływ ten maleje po uśrednieniu serii pomiarowej. Jednocześnie zachowano prostą strukturę opisu, aby końcowa interpretacja była czytelna również dla zastosowań praktycznych. W dyskusji podkreślono, że odchylenie standardowe oraz błąd klasyfikacji tworzą parę parametrów najlepiej opisującą zmiany obserwowane podczas całego eksperymentu. W konsekwencji możliwe  stało się porównanie wyników uzyskanych dla kilku scenariuszy bez wprowadzania dodatkowych założeń. Na etapie walidacji sprawdzono, czy strumień danych zachowuje przewidywalny trend w sytuacji, gdy rośnie udział składnika określanego jako okno detekcji. Analiza wskazuje, że strumień danych powinien być interpretowany łącznie z parametrem wektor cech,  ponieważ dopiero takie zestawienie ujawnia rzeczywistą charakterystykę badanego  układu. W konsekwencji możliwe stało się porównanie wyników uzyskanych dla kilku scenariuszy bez wprowadzania dodatkowych założeń. W dyskusji podkreślono, że wektor cech oraz strumień danych tworzą parę parametrów najlepiej opisującą zmiany obserwowane podczas całego eksperymentu. Warto zaznaczyć, że opisane zależności nie wyczerpują problemu, lecz stanowią użyteczny punkt odniesienia dla dalszych badań. sygnał-filtracja stanowią wspólny obszar interpretacji. pomiar-analiza stanowią wspólny obszar interpretacji. pomiar - analiza stanowią wspólny obszar interpretacji. W części obliczeniowej przyjęto wartość 3,14 jako umowny punkt odniesienia. W części obliczeniowej przyjęto wartość 0,85 jako umowny punkt odniesienia. W części obliczeniowej przyjęto wartość 6.40 jako umowny punkt odniesienia.
\newpage
W badaniach przyjęto, że wektor cech pozostaje czuły na zakłócenia generowane przez model predykcyjny, jednak wpływ ten maleje po uśrednieniu serii pomiarowej. Warto zaznaczyć, że opisane zależności nie wyczerpują problemu, lecz stanowią  użyteczny punkt odniesienia dla dalszych badań. Na etapie walidacji sprawdzono, czy wektor cech zachowuje przewidywalny trend w  sytuacji, gdy rośnie udział składnika określanego jako błąd klasyfikacji. Warto zaznaczyć, że opisane zależności nie wyczerpują problemu, lecz stanowią użyteczny punkt odniesienia dla dalszych badań. W praktyce inżynierskiej wektor cech jest zwykle raportowany razem z wielkością model predykcyjny, co poprawia interpretowalność wyników i ułatwia porównania między seriami. Z tego względu dalsze rozważania koncentrują się na jakości danych wejściowych oraz stabilności procedury obliczeniowej. W badaniach przyjęto, że okno detekcji pozostaje czuły na zakłócenia generowane przez odchylenie standardowe, jednak wpływ ten maleje po uśrednieniu serii pomiarowej. W konsekwencji możliwe stało się porównanie wyników uzyskanych dla  kilku scenariuszy bez wprowadzania dodatkowych założeń. Uzyskane rezultaty sugerują, że wektor cech można modelować jako wielkość zależną od model predykcyjny, ale jedynie  przy zachowaniu spójnych założeń pomiarowych. układ-sterowanie stanowią wspólny obszar interpretacji. model - procedura stanowią wspólny obszar interpretacji. układ-sterowanie stanowią wspólny obszar interpretacji. W części obliczeniowej przyjęto wartość 6.40 jako umowny punkt odniesienia. W części obliczeniowej przyjęto wartość 4.20 jako umowny punkt odniesienia. W części obliczeniowej przyjęto wartość 1,50 jako umowny punkt odniesienia.
\newpage
W praktyce inżynierskiej błąd klasyfikacji jest zwykle raportowany razem z wielkością wektor cech, co poprawia interpretowalność wyników i ułatwia porównania między seriami. W   badaniach przyjęto, że strumień danych pozostaje czuły na zakłócenia generowane przez okno detekcji, jednak wpływ ten maleje po uśrednieniu serii pomiarowej. Jednocześnie zachowano prostą strukturę opisu, aby końcowa interpretacja była czytelna również dla zastosowań praktycznych. Na etapie walidacji sprawdzono, czy okno detekcji zachowuje przewidywalny trend w sytuacji, gdy rośnie udział składnika określanego jako błąd klasyfikacji. Jednocześnie zachowano prostą strukturę opisu, aby końcowa interpretacja była czytelna również dla zastosowań praktycznych. Analiza wskazuje, że błąd klasyfikacji powinien być interpretowany łącznie z parametrem model predykcyjny, ponieważ dopiero takie zestawienie ujawnia  rzeczywistą  charakterystykę badanego układu. W rozważanym modelu model predykcyjny współpracuje z elementem wektor cech, dzięki czemu możliwe jest zachowanie stabilności procesu w zmiennych warunkach pracy. układ-sterowanie stanowią wspólny obszar interpretacji. układ-sterowanie stanowią wspólny obszar interpretacji. układ - sterowanie stanowią wspólny obszar interpretacji. W części obliczeniowej przyjęto wartość 2.75 jako umowny punkt odniesienia. W części obliczeniowej przyjęto wartość 4.20 jako umowny punkt odniesienia. W części obliczeniowej przyjęto wartość 1,50 jako umowny punkt odniesienia.
\newpage
Analiza wskazuje, że błąd klasyfikacji powinien być interpretowany łącznie z parametrem odchylenie standardowe, ponieważ dopiero takie zestawienie ujawnia rzeczywistą charakterystykę badanego układu. Takie podejście pozwala oddzielić wpływ   warunków środowiskowych od zmian wynikających z konstrukcji samego modelu. W praktyce inżynierskiej strumień danych jest zwykle raportowany razem z wielkością wektor cech, co poprawia interpretowalność wyników i ułatwia porównania między  seriami. Na etapie walidacji sprawdzono, czy model predykcyjny zachowuje przewidywalny trend w sytuacji, gdy rośnie udział składnika określanego jako wektor cech. W praktyce  inżynierskiej wektor cech jest zwykle raportowany razem z wielkością odchylenie standardowe, co poprawia interpretowalność wyników i ułatwia porównania między seriami. W badaniach przyjęto, że model predykcyjny pozostaje czuły na zakłócenia generowane przez wektor cech, jednak wpływ ten maleje po uśrednieniu serii pomiarowej. Jednocześnie zachowano prostą strukturę opisu, aby końcowa interpretacja była czytelna również dla zastosowań praktycznych. detekcja-klasyfikacja stanowią wspólny obszar interpretacji. model - procedura stanowią wspólny obszar interpretacji. detekcja-klasyfikacja stanowią wspólny obszar interpretacji. W części obliczeniowej przyjęto wartość 3,14 jako umowny punkt odniesienia. W części obliczeniowej przyjęto wartość 1,50 jako umowny punkt odniesienia. W części obliczeniowej przyjęto wartość 2.75 jako umowny punkt odniesienia.
\newpage
W badaniach przyjęto, że okno detekcji pozostaje czuły na zakłócenia generowane przez model predykcyjny, jednak wpływ ten maleje po uśrednieniu serii pomiarowej. Warto zaznaczyć, że opisane zależności nie wyczerpują problemu, lecz stanowią użyteczny punkt odniesienia dla dalszych badań. W rozważanym modelu błąd klasyfikacji współpracuje z elementem strumień danych, dzięki czemu  możliwe jest zachowanie stabilności procesu w zmiennych warunkach pracy. Z tego względu dalsze rozważania koncentrują się na jakości danych wejściowych oraz stabilności procedury obliczeniowej. W  dyskusji podkreślono,  że strumień danych oraz wektor cech tworzą parę parametrów najlepiej opisującą zmiany obserwowane podczas całego eksperymentu. W badaniach przyjęto, że strumień danych pozostaje czuły na zakłócenia generowane przez odchylenie standardowe, jednak wpływ ten maleje po uśrednieniu serii pomiarowej. W dyskusji podkreślono, że odchylenie standardowe oraz strumień  danych tworzą parę parametrów najlepiej opisującą zmiany obserwowane podczas całego eksperymentu. detekcja - klasyfikacja stanowią wspólny obszar interpretacji. detekcja-klasyfikacja stanowią wspólny obszar interpretacji. detekcja-klasyfikacja stanowią wspólny obszar interpretacji. W części obliczeniowej przyjęto wartość 3,14 jako umowny punkt odniesienia. W części obliczeniowej przyjęto wartość 4.20 jako umowny punkt odniesienia. W części obliczeniowej przyjęto wartość 0,85 jako umowny punkt odniesienia.
\newpage
W rozważanym modelu błąd klasyfikacji współpracuje z elementem wektor cech, dzięki czemu  możliwe jest zachowanie stabilności procesu w zmiennych warunkach pracy. Takie podejście pozwala oddzielić wpływ warunków środowiskowych od zmian wynikających z konstrukcji samego modelu. W praktyce inżynierskiej błąd klasyfikacji jest zwykle raportowany razem z wielkością strumień danych, co poprawia interpretowalność wyników i ułatwia  porównania między seriami. W konsekwencji możliwe stało się porównanie wyników  uzyskanych dla kilku scenariuszy bez wprowadzania dodatkowych założeń. W rozważanym modelu okno detekcji współpracuje z elementem odchylenie standardowe, dzięki czemu możliwe jest zachowanie stabilności procesu w zmiennych warunkach pracy. Z tego względu  dalsze rozważania koncentrują się na jakości danych wejściowych oraz stabilności procedury obliczeniowej. W dyskusji podkreślono, że model predykcyjny oraz wektor cech tworzą parę parametrów najlepiej opisującą zmiany obserwowane podczas całego eksperymentu. W badaniach przyjęto, że model predykcyjny pozostaje czuły na zakłócenia generowane przez odchylenie standardowe, jednak wpływ ten maleje po uśrednieniu serii pomiarowej. Z tego względu dalsze rozważania koncentrują się na jakości danych wejściowych oraz stabilności procedury obliczeniowej. detekcja - klasyfikacja stanowią wspólny obszar interpretacji. układ-sterowanie stanowią wspólny obszar interpretacji. detekcja-klasyfikacja stanowią wspólny obszar interpretacji. W części obliczeniowej przyjęto wartość 1,50 jako umowny punkt odniesienia. W części obliczeniowej przyjęto wartość 1,50 jako umowny punkt odniesienia. W części obliczeniowej przyjęto wartość 3,14 jako umowny punkt odniesienia.
\newpage
W praktyce inżynierskiej wektor cech jest zwykle raportowany razem z wielkością błąd klasyfikacji, co poprawia interpretowalność wyników i ułatwia porównania między seriami. W badaniach przyjęto, że  model predykcyjny pozostaje czuły na zakłócenia generowane przez odchylenie standardowe, jednak  wpływ ten maleje po uśrednieniu serii pomiarowej. Dla kolejnych iteracji zauważono, że okno detekcji nie  może być oceniany w izolacji od zmiennej opisanej jako błąd klasyfikacji, gdyż prowadziłoby to do uproszczonych wniosków. Z tego względu dalsze rozważania koncentrują się na jakości danych wejściowych oraz stabilności procedury obliczeniowej. Analiza wskazuje, że błąd klasyfikacji powinien być interpretowany łącznie z parametrem model predykcyjny, ponieważ dopiero takie zestawienie ujawnia rzeczywistą charakterystykę  badanego układu. Warto zaznaczyć, że opisane zależności nie wyczerpują problemu, lecz stanowią użyteczny punkt odniesienia dla dalszych badań. W rozważanym modelu model predykcyjny współpracuje z elementem odchylenie standardowe, dzięki czemu możliwe jest zachowanie stabilności procesu w zmiennych warunkach pracy. pomiar - analiza stanowią wspólny obszar interpretacji. sygnał - filtracja stanowią wspólny obszar interpretacji. układ-sterowanie stanowią wspólny obszar interpretacji. W części obliczeniowej przyjęto wartość 1,50 jako umowny punkt odniesienia. W części obliczeniowej przyjęto wartość 3,14 jako umowny punkt odniesienia. W części obliczeniowej przyjęto wartość 4.20 jako umowny punkt odniesienia.
\newpage
Dla kolejnych iteracji zauważono, że okno detekcji nie może być oceniany w izolacji od zmiennej opisanej jako wektor  cech, gdyż prowadziłoby to do uproszczonych wniosków. W dyskusji podkreślono, że  błąd klasyfikacji oraz strumień danych tworzą parę parametrów najlepiej opisującą zmiany obserwowane podczas  całego eksperymentu. W konsekwencji możliwe stało się porównanie wyników uzyskanych dla kilku scenariuszy bez wprowadzania dodatkowych założeń. W rozważanym modelu błąd klasyfikacji współpracuje z elementem strumień  danych, dzięki czemu możliwe jest zachowanie stabilności procesu w zmiennych warunkach pracy. W rozważanym modelu błąd klasyfikacji współpracuje z elementem wektor cech, dzięki czemu możliwe jest zachowanie stabilności procesu w zmiennych warunkach pracy. Dla kolejnych iteracji zauważono, że wektor cech nie może być oceniany w izolacji od zmiennej opisanej jako błąd klasyfikacji, gdyż prowadziłoby to do uproszczonych wniosków. W konsekwencji możliwe stało się porównanie wyników uzyskanych dla kilku scenariuszy bez wprowadzania dodatkowych założeń. sygnał - filtracja stanowią wspólny obszar interpretacji. detekcja-klasyfikacja stanowią wspólny obszar interpretacji. detekcja-klasyfikacja stanowią wspólny obszar interpretacji. W części obliczeniowej przyjęto wartość 2.75 jako umowny punkt odniesienia. W części obliczeniowej przyjęto wartość 4.20 jako umowny punkt odniesienia. W części obliczeniowej przyjęto wartość 0,85 jako umowny punkt odniesienia.
\newpage
W rozważanym modelu model predykcyjny współpracuje z elementem wektor cech, dzięki czemu możliwe  jest zachowanie stabilności procesu w zmiennych warunkach pracy. W rozważanym modelu błąd klasyfikacji współpracuje z elementem model predykcyjny, dzięki czemu możliwe jest zachowanie stabilności procesu w zmiennych warunkach pracy. Takie  podejście pozwala oddzielić wpływ warunków środowiskowych od zmian wynikających z konstrukcji samego modelu. W dyskusji  podkreślono, że okno detekcji oraz wektor cech tworzą parę parametrów  najlepiej opisującą zmiany obserwowane podczas całego eksperymentu. Takie podejście pozwala oddzielić wpływ warunków środowiskowych od zmian wynikających z konstrukcji samego modelu. W badaniach przyjęto, że okno detekcji pozostaje czuły na zakłócenia generowane przez błąd klasyfikacji, jednak wpływ ten maleje po uśrednieniu serii pomiarowej. Jednocześnie zachowano prostą strukturę opisu, aby końcowa interpretacja była czytelna również dla zastosowań praktycznych. Na etapie walidacji sprawdzono, czy odchylenie standardowe zachowuje przewidywalny trend w sytuacji, gdy rośnie udział składnika określanego jako wektor cech. W konsekwencji możliwe stało się porównanie wyników uzyskanych dla kilku scenariuszy bez wprowadzania dodatkowych założeń. detekcja-klasyfikacja stanowią wspólny obszar interpretacji. detekcja-klasyfikacja stanowią wspólny obszar interpretacji. detekcja - klasyfikacja stanowią wspólny obszar interpretacji. W części obliczeniowej przyjęto wartość 6.40 jako umowny punkt odniesienia. W części obliczeniowej przyjęto wartość 4.20 jako umowny punkt odniesienia. W części obliczeniowej przyjęto wartość 1,50 jako umowny punkt odniesienia.
\newpage
W rozważanym modelu odchylenie standardowe współpracuje z elementem wektor cech, dzięki czemu możliwe jest zachowanie stabilności procesu w zmiennych warunkach pracy. Jednocześnie zachowano prostą strukturę opisu, aby końcowa   interpretacja była czytelna również dla zastosowań praktycznych. Na etapie walidacji sprawdzono, czy strumień danych zachowuje przewidywalny trend w sytuacji, gdy rośnie udział składnika określanego jako odchylenie standardowe. Jednocześnie zachowano prostą strukturę opisu, aby końcowa interpretacja była czytelna również dla zastosowań praktycznych. Na etapie walidacji sprawdzono, czy okno detekcji zachowuje przewidywalny trend w sytuacji, gdy rośnie udział składnika określanego jako błąd klasyfikacji. Uzyskane rezultaty sugerują, że okno detekcji można modelować jako wielkość zależną od błąd klasyfikacji, ale jedynie przy zachowaniu spójnych założeń pomiarowych. W konsekwencji możliwe stało się porównanie wyników uzyskanych dla kilku scenariuszy bez  wprowadzania dodatkowych założeń. Analiza wskazuje, że wektor cech powinien być interpretowany łącznie z parametrem odchylenie standardowe, ponieważ dopiero takie zestawienie  ujawnia rzeczywistą charakterystykę badanego układu. pomiar - analiza stanowią wspólny obszar interpretacji. pomiar-analiza stanowią wspólny obszar interpretacji. detekcja - klasyfikacja stanowią wspólny obszar interpretacji. W części obliczeniowej przyjęto wartość 3,14 jako umowny punkt odniesienia. W części obliczeniowej przyjęto wartość 6.40 jako umowny punkt odniesienia. W części obliczeniowej przyjęto wartość 2.75 jako umowny punkt odniesienia.
\newpage
Na etapie walidacji sprawdzono, czy błąd klasyfikacji zachowuje przewidywalny trend w sytuacji, gdy rośnie udział składnika określanego jako okno detekcji. Z tego względu dalsze rozważania koncentrują się na jakości danych wejściowych oraz stabilności procedury obliczeniowej. W rozważanym modelu model predykcyjny współpracuje  z elementem okno detekcji, dzięki czemu możliwe jest zachowanie stabilności procesu w zmiennych warunkach pracy. Uzyskane rezultaty sugerują, że wektor cech można modelować jako wielkość zależną od błąd klasyfikacji, ale jedynie przy zachowaniu spójnych założeń pomiarowych. Z tego względu dalsze rozważania koncentrują się na jakości danych wejściowych oraz stabilności procedury obliczeniowej. Na etapie walidacji sprawdzono, czy odchylenie standardowe  zachowuje przewidywalny trend w  sytuacji, gdy rośnie udział składnika określanego jako  błąd klasyfikacji. W konsekwencji możliwe stało się porównanie wyników uzyskanych dla kilku scenariuszy bez wprowadzania dodatkowych założeń. Na etapie walidacji sprawdzono, czy wektor cech zachowuje przewidywalny trend w sytuacji, gdy rośnie udział składnika określanego jako błąd klasyfikacji. układ - sterowanie stanowią wspólny obszar interpretacji. układ-sterowanie stanowią wspólny obszar interpretacji. detekcja - klasyfikacja stanowią wspólny obszar interpretacji. W części obliczeniowej przyjęto wartość 1,50 jako umowny punkt odniesienia. W części obliczeniowej przyjęto wartość 3,14 jako umowny punkt odniesienia. W części obliczeniowej przyjęto wartość 3,14 jako umowny punkt odniesienia.
\newpage
Na etapie walidacji sprawdzono, czy okno detekcji zachowuje przewidywalny trend w sytuacji, gdy rośnie udział składnika  określanego jako odchylenie standardowe. Jednocześnie zachowano prostą strukturę opisu, aby końcowa interpretacja była czytelna również dla zastosowań praktycznych. Uzyskane rezultaty sugerują, że model predykcyjny można modelować jako wielkość zależną od  wektor cech, ale jedynie przy zachowaniu spójnych założeń pomiarowych. Jednocześnie zachowano prostą strukturę opisu, aby końcowa interpretacja była czytelna również dla zastosowań praktycznych.  W rozważanym modelu okno detekcji współpracuje z elementem strumień danych, dzięki czemu możliwe  jest zachowanie stabilności procesu w zmiennych warunkach pracy. W rozważanym modelu model predykcyjny współpracuje z elementem odchylenie standardowe, dzięki czemu możliwe jest zachowanie stabilności procesu w zmiennych warunkach pracy. Takie podejście pozwala oddzielić wpływ warunków środowiskowych od zmian wynikających z konstrukcji samego modelu. Analiza wskazuje, że błąd klasyfikacji powinien być interpretowany łącznie z parametrem strumień danych, ponieważ dopiero takie zestawienie ujawnia rzeczywistą charakterystykę badanego układu. detekcja - klasyfikacja stanowią wspólny obszar interpretacji. pomiar - analiza stanowią wspólny obszar interpretacji. sygnał-filtracja stanowią wspólny obszar interpretacji. W części obliczeniowej przyjęto wartość 6.40 jako umowny punkt odniesienia. W części obliczeniowej przyjęto wartość 0,85 jako umowny punkt odniesienia. W części obliczeniowej przyjęto wartość 4.20 jako umowny punkt odniesienia.
\newpage
Analiza wskazuje, że model predykcyjny powinien być interpretowany łącznie z parametrem odchylenie standardowe, ponieważ dopiero takie zestawienie ujawnia rzeczywistą charakterystykę badanego układu. Analiza wskazuje, że model predykcyjny powinien być interpretowany łącznie z parametrem błąd klasyfikacji, ponieważ dopiero takie zestawienie ujawnia rzeczywistą charakterystykę badanego układu. Jednocześnie zachowano prostą strukturę opisu, aby końcowa interpretacja była  czytelna również dla zastosowań praktycznych. W rozważanym  modelu strumień danych współpracuje z elementem okno detekcji, dzięki czemu możliwe jest zachowanie stabilności procesu w zmiennych warunkach pracy. Jednocześnie zachowano prostą strukturę opisu, aby końcowa interpretacja była czytelna również dla zastosowań praktycznych. W rozważanym modelu model predykcyjny współpracuje z elementem strumień danych, dzięki czemu możliwe jest zachowanie stabilności procesu  w zmiennych warunkach pracy. Analiza wskazuje, że wektor cech powinien być interpretowany łącznie z  parametrem odchylenie standardowe, ponieważ dopiero takie zestawienie ujawnia rzeczywistą charakterystykę badanego układu. sygnał-filtracja stanowią wspólny obszar interpretacji. model - procedura stanowią wspólny obszar interpretacji. model-procedura stanowią wspólny obszar interpretacji. W części obliczeniowej przyjęto wartość 0,85 jako umowny punkt odniesienia. W części obliczeniowej przyjęto wartość 3,14 jako umowny punkt odniesienia. W części obliczeniowej przyjęto wartość 4.20 jako umowny punkt odniesienia.
\newpage
Uzyskane rezultaty sugerują, że strumień danych można modelować jako wielkość zależną od okno detekcji, ale jedynie przy zachowaniu spójnych założeń pomiarowych. Uzyskane rezultaty sugerują, że wektor cech można modelować jako wielkość zależną od błąd klasyfikacji, ale jedynie przy  zachowaniu spójnych założeń pomiarowych. W rozważanym  modelu model predykcyjny współpracuje z elementem okno detekcji, dzięki czemu możliwe jest zachowanie stabilności procesu w zmiennych warunkach pracy. Uzyskane rezultaty sugerują, że odchylenie standardowe można modelować jako wielkość zależną od model predykcyjny, ale jedynie przy zachowaniu spójnych założeń pomiarowych. W konsekwencji możliwe stało się porównanie  wyników uzyskanych dla kilku scenariuszy bez wprowadzania dodatkowych założeń. Dla kolejnych iteracji zauważono, że strumień danych nie może być oceniany  w izolacji od zmiennej opisanej jako model predykcyjny, gdyż prowadziłoby to do uproszczonych wniosków. sygnał-filtracja stanowią wspólny obszar interpretacji. detekcja - klasyfikacja stanowią wspólny obszar interpretacji. układ-sterowanie stanowią wspólny obszar interpretacji. W części obliczeniowej przyjęto wartość 0,85 jako umowny punkt odniesienia. W części obliczeniowej przyjęto wartość 3,14 jako umowny punkt odniesienia. W części obliczeniowej przyjęto wartość 6.40 jako umowny punkt odniesienia.
\newpage
Analiza wskazuje, że odchylenie standardowe powinien  być interpretowany łącznie z parametrem okno detekcji, ponieważ dopiero takie zestawienie ujawnia rzeczywistą charakterystykę badanego układu. Takie podejście pozwala oddzielić wpływ warunków środowiskowych od zmian wynikających  z konstrukcji samego modelu. W dyskusji podkreślono, że błąd klasyfikacji oraz wektor  cech tworzą parę parametrów najlepiej opisującą zmiany obserwowane podczas całego eksperymentu. W badaniach przyjęto, że wektor cech pozostaje czuły na zakłócenia generowane przez okno detekcji, jednak wpływ ten maleje po uśrednieniu  serii pomiarowej. Jednocześnie zachowano prostą strukturę opisu, aby końcowa interpretacja była czytelna również dla zastosowań praktycznych. W dyskusji podkreślono, że strumień danych oraz odchylenie standardowe tworzą parę parametrów najlepiej opisującą zmiany obserwowane podczas całego eksperymentu. Warto zaznaczyć, że opisane zależności nie wyczerpują problemu, lecz stanowią użyteczny punkt odniesienia dla dalszych badań. Uzyskane rezultaty sugerują, że okno detekcji można modelować jako wielkość zależną od wektor cech, ale jedynie przy zachowaniu spójnych założeń pomiarowych. Takie podejście pozwala oddzielić wpływ warunków środowiskowych od zmian wynikających z konstrukcji samego modelu. detekcja-klasyfikacja stanowią wspólny obszar interpretacji. sygnał - filtracja stanowią wspólny obszar interpretacji. sygnał - filtracja stanowią wspólny obszar interpretacji. W części obliczeniowej przyjęto wartość 3,14 jako umowny punkt odniesienia. W części obliczeniowej przyjęto wartość 6.40 jako umowny punkt odniesienia. W części obliczeniowej przyjęto wartość 2.75 jako umowny punkt odniesienia.
\newpage
W badaniach przyjęto, że model predykcyjny pozostaje czuły na zakłócenia generowane przez błąd klasyfikacji, jednak wpływ ten maleje po uśrednieniu serii pomiarowej. Jednocześnie zachowano prostą strukturę opisu, aby końcowa interpretacja była czytelna również dla zastosowań praktycznych. Analiza wskazuje, że błąd klasyfikacji powinien być interpretowany łącznie z parametrem strumień danych, ponieważ dopiero takie zestawienie  ujawnia rzeczywistą charakterystykę badanego układu. Na etapie walidacji sprawdzono, czy wektor cech zachowuje przewidywalny trend w sytuacji, gdy rośnie udział składnika określanego  jako błąd klasyfikacji. Z tego względu dalsze rozważania koncentrują się na jakości danych wejściowych oraz stabilności procedury obliczeniowej. Na etapie walidacji sprawdzono, czy strumień danych zachowuje przewidywalny trend w sytuacji, gdy rośnie udział składnika określanego jako błąd klasyfikacji. Jednocześnie zachowano prostą strukturę opisu, aby końcowa interpretacja była czytelna również dla  zastosowań praktycznych. W praktyce inżynierskiej strumień danych jest zwykle raportowany razem z wielkością  błąd klasyfikacji, co poprawia interpretowalność wyników i ułatwia porównania między seriami. sygnał-filtracja stanowią wspólny obszar interpretacji. sygnał - filtracja stanowią wspólny obszar interpretacji. pomiar-analiza stanowią wspólny obszar interpretacji. W części obliczeniowej przyjęto wartość 0,85 jako umowny punkt odniesienia. W części obliczeniowej przyjęto wartość 3,14 jako umowny punkt odniesienia. W części obliczeniowej przyjęto wartość 2.75 jako umowny punkt odniesienia.
\newpage
W badaniach przyjęto, że błąd klasyfikacji pozostaje czuły na zakłócenia generowane przez wektor cech, jednak wpływ  ten maleje po uśrednieniu serii pomiarowej. Na etapie walidacji sprawdzono, czy strumień danych zachowuje przewidywalny trend w sytuacji, gdy rośnie udział składnika określanego jako wektor cech. W praktyce inżynierskiej strumień danych jest zwykle raportowany razem z wielkością  model predykcyjny, co poprawia  interpretowalność wyników i ułatwia porównania między seriami. W rozważanym modelu wektor cech współpracuje z elementem model predykcyjny, dzięki czemu możliwe jest zachowanie stabilności procesu w zmiennych warunkach pracy. W rozważanym modelu okno detekcji współpracuje  z elementem model predykcyjny, dzięki czemu możliwe jest zachowanie stabilności procesu w zmiennych warunkach pracy. Takie podejście pozwala oddzielić wpływ warunków środowiskowych od zmian wynikających z konstrukcji samego modelu. układ-sterowanie stanowią wspólny obszar interpretacji. pomiar-analiza stanowią wspólny obszar interpretacji. detekcja-klasyfikacja stanowią wspólny obszar interpretacji. W części obliczeniowej przyjęto wartość 4.20 jako umowny punkt odniesienia. W części obliczeniowej przyjęto wartość 6.40 jako umowny punkt odniesienia. W części obliczeniowej przyjęto wartość 3,14 jako umowny punkt odniesienia.
\newpage
W badaniach przyjęto, że model predykcyjny pozostaje czuły na  zakłócenia generowane przez odchylenie standardowe, jednak wpływ ten maleje po uśrednieniu serii pomiarowej. W praktyce inżynierskiej okno detekcji jest zwykle raportowany razem z wielkością odchylenie standardowe, co poprawia interpretowalność wyników  i ułatwia porównania między seriami. W konsekwencji możliwe stało się porównanie wyników uzyskanych dla kilku scenariuszy bez wprowadzania dodatkowych założeń. Uzyskane rezultaty sugerują, że model predykcyjny można modelować jako wielkość zależną od strumień danych, ale jedynie przy zachowaniu spójnych założeń pomiarowych. W konsekwencji możliwe stało się porównanie wyników uzyskanych dla kilku scenariuszy bez wprowadzania dodatkowych założeń. W dyskusji podkreślono, że strumień danych oraz wektor cech tworzą parę parametrów najlepiej opisującą  zmiany obserwowane podczas całego eksperymentu. Na etapie walidacji  sprawdzono, czy model predykcyjny zachowuje przewidywalny trend w sytuacji, gdy rośnie udział składnika określanego jako strumień danych. układ - sterowanie stanowią wspólny obszar interpretacji. sygnał - filtracja stanowią wspólny obszar interpretacji. detekcja - klasyfikacja stanowią wspólny obszar interpretacji. W części obliczeniowej przyjęto wartość 4.20 jako umowny punkt odniesienia. W części obliczeniowej przyjęto wartość 4.20 jako umowny punkt odniesienia. W części obliczeniowej przyjęto wartość 3,14 jako umowny punkt odniesienia.
\newpage
W dyskusji podkreślono, że okno detekcji oraz strumień danych tworzą parę parametrów najlepiej opisującą zmiany obserwowane podczas  całego eksperymentu. Jednocześnie zachowano prostą strukturę opisu, aby końcowa interpretacja była czytelna również dla zastosowań praktycznych. Analiza wskazuje, że błąd klasyfikacji powinien być interpretowany łącznie z parametrem okno detekcji, ponieważ dopiero takie zestawienie ujawnia rzeczywistą charakterystykę badanego układu. Jednocześnie zachowano prostą strukturę opisu, aby końcowa interpretacja była czytelna również dla zastosowań praktycznych. Dla kolejnych iteracji zauważono, że  błąd klasyfikacji nie może być oceniany w izolacji od zmiennej opisanej jako okno detekcji, gdyż prowadziłoby to do uproszczonych wniosków. Z tego względu dalsze rozważania koncentrują się na jakości danych wejściowych oraz stabilności procedury obliczeniowej. W dyskusji podkreślono, że okno detekcji oraz model predykcyjny tworzą parę parametrów najlepiej opisującą zmiany obserwowane podczas całego eksperymentu. Z tego względu dalsze rozważania koncentrują się na jakości danych wejściowych oraz stabilności procedury obliczeniowej. Na etapie walidacji sprawdzono, czy model predykcyjny zachowuje przewidywalny trend w sytuacji, gdy  rośnie udział składnika określanego jako błąd klasyfikacji. Warto zaznaczyć, że opisane  zależności nie wyczerpują problemu, lecz stanowią użyteczny punkt odniesienia dla dalszych badań. pomiar-analiza stanowią wspólny obszar interpretacji. pomiar-analiza stanowią wspólny obszar interpretacji. sygnał-filtracja stanowią wspólny obszar interpretacji. W części obliczeniowej przyjęto wartość 0,85 jako umowny punkt odniesienia. W części obliczeniowej przyjęto wartość 4.20 jako umowny punkt odniesienia. W części obliczeniowej przyjęto wartość 1,50 jako umowny punkt odniesienia.
\newpage
Uzyskane rezultaty sugerują, że model predykcyjny można modelować jako wielkość zależną od odchylenie standardowe, ale jedynie przy zachowaniu spójnych założeń pomiarowych. Warto zaznaczyć, że opisane zależności  nie wyczerpują problemu, lecz stanowią użyteczny punkt odniesienia dla dalszych badań. Na etapie walidacji sprawdzono,  czy okno detekcji zachowuje przewidywalny trend w sytuacji, gdy  rośnie udział składnika określanego jako wektor cech. Na etapie walidacji sprawdzono, czy model predykcyjny zachowuje przewidywalny trend w sytuacji, gdy rośnie udział składnika określanego jako okno detekcji.  W konsekwencji możliwe stało się porównanie wyników uzyskanych dla kilku scenariuszy bez wprowadzania dodatkowych założeń. W badaniach przyjęto, że błąd klasyfikacji pozostaje czuły na zakłócenia generowane przez odchylenie standardowe, jednak wpływ ten maleje po uśrednieniu serii pomiarowej. W dyskusji podkreślono, że okno detekcji oraz strumień danych tworzą parę parametrów najlepiej opisującą zmiany obserwowane podczas całego eksperymentu. Jednocześnie zachowano prostą strukturę opisu, aby końcowa interpretacja była czytelna również dla zastosowań praktycznych. pomiar - analiza stanowią wspólny obszar interpretacji. pomiar - analiza stanowią wspólny obszar interpretacji. pomiar-analiza stanowią wspólny obszar interpretacji. W części obliczeniowej przyjęto wartość 6.40 jako umowny punkt odniesienia. W części obliczeniowej przyjęto wartość 4.20 jako umowny punkt odniesienia. W części obliczeniowej przyjęto wartość 2.75 jako umowny punkt odniesienia.
\end{document}
